\subsection{Gluing predictable compensators and bounding jumps}
For each $n$, let $P^n$ be the dual predictable projection (compensator) of $J^{\tau_n}$ and
$Q^n=J^{\tau_n}-P^n$ its martingale residual. The local construction gives
ordered overlaps
\begin{equation}
 m\leq n\quad\Longrightarrow\quad (P^n)^{\tau_m}\sim P^m.
 \label{eq:compensator-ordered}
\end{equation}
The gluing theorem for locally finite-variation processes instead requires
\begin{equation}
 (P^n)^{\tau_n\wedge\tau_m}\sim
 (P^m)^{\tau_n\wedge\tau_m}\qquad(n,m\in\mathbb N).
 \label{eq:compensator-minimum}
\end{equation}
Taking $m=n$ in \eqref{eq:compensator-ordered} shows that each coordinate is unchanged by its own
stop. On one full-measure set the sequence $(\tau_n)$ is increasing; hence
$\tau_n\wedge\tau_m=\tau_m$ when $m\leq n$.
Stopping absorption and the total order on the indices prove \eqref{eq:compensator-minimum}.
Thus the ordered overlaps supply the all-pairs compatibility needed for gluing.

The local limit in \eqref{eq:compensator-minimum} is eventually constant at each point, and
predictability passes to that limit. Compatibility and the almost-sure
monotonicity and exhaustion of the schedule give right continuity and
local finite variation almost surely. Set the whole path to zero on the
exceptional null set, obtaining a predictable version $P$ with these path
properties everywhere. No replacement of the original compatibility or
schedule by pointwise monotone data is necessary.
Initially, however, $P^n_0=0$ and therefore $P_0=0$ hold only almost surely.
The exceptional set $\{P_0\ne0\}$ is measurable at time zero. Set the whole
path to zero there too. This second normalization preserves predictability,
right continuity, local finite variation and stopped indistinguishability
with every coordinate. All these properties now hold for the same version $P$.

The closed sources have the same minimum-stop compatibility as \eqref{eq:compensator-minimum};
subtracting the projections gives it for $(Q^n)$ too. These coordinates
are c\`adl\`ag true martingales. Martingale gluing produces one strongly
adapted c\`adl\`ag local martingale $\widetilde Q$, whose stop at every
$\tau_n$ is indistinguishable from $Q^n$. Thus
\[
 J^{\tau_n}\sim\widetilde Q^{\tau_n}+P^{\tau_n}.
\]
Intersect the countably many stopped-agreement events with the exhaustion
event. For any $t$, some $n$ satisfies $t\leq\tau_n$, proving
$J\sim\widetilde Q+P$. On the normalized representatives define $Q=J-P$.
Then $Q\sim\widetilde Q$, so $Q$ is a local martingale; it is also adapted,
c\`adl\`ag, locally of finite variation, and zero-start.
This proof uses gluing of true-martingale coordinates and transfer by
indistinguishability, rather than expanding the time-zero indicators in
the definition of a local property.

Set $L=N-Q=R+P$. Then $N=L+Q$, and $L$ is a c\`adl\`ag local martingale.
In addition to $|\Delta R|\leq c$, we need the following compensator estimate.
\begin{lemma}[Compensator jump estimate]\label{lem:compensator-jump-bound}
For fixed $c>0$ and $T<\infty$, on one full-measure set,
\[
 0<t\leq T\quad\Longrightarrow\quad
 |\Delta P_t|\leq c,\qquad |\Delta L_t|\leq2c.
\]
\end{lemma}
\begin{proof}
Stop $R=L+(-P)$ at $T$. All nonzero jumps of the predictable c\`adl\`ag
finite-variation process $P^T$ are covered by countably many announced
predictable stopping graphs $(\sigma_k)$, constructed from rational level
crossings of cumulative variation. Each completed time is at most $T+1$.
Apply the localized conditional-expectation estimate to $L^T$ and $-P^T$.
Without assuming integrability of all jumps of the original local
martingale, this gives
\[
 |\Delta(P^T)_{\sigma_k}|
 \leq E[c\mid\F_{\sigma_k-}]=c\quad\text{a.s.}
\]
Its input is $|\Delta(R^T)_t|\leq c$; after $T$ the stopped process has
zero jumps. Combine the countably many null sets and use the graph
coverage. Times with zero jump of $P$ satisfy the estimate automatically.
This proves the first bound. The second follows from
$\Delta L=\Delta R+\Delta P$ and the triangle inequality.
For $T=0$ the relevant interval is empty. At the dummy enumeration value
$T+1$, the stopped process has zero jump.
\end{proof}
Thus compensation first gives the bound $c$ for $\Delta P$, and then the
bound $2c$ for $\Delta L$; it does not assert a bound $c$ for $\Delta Q$.
From the original zero-start c\`adl\`ag local martingale, a threshold,
a finite horizon, and the usual conditions, the construction supplies
$N=L+Q$ with these same components. Both components are zero-start
c\`adl\`ag local martingales and $Q$ is adapted and locally of finite
variation. Predictability of $Q$ is not asserted.
For every pathwise stop $\rho\leq T$, the bound
$|\Delta(L^\rho)_t|\leq2c$ also holds at every time outside the same null
set. The decomposition at this stage depends on the fixed horizon.

\begin{lemma}[Compatibility between horizons]
Write $P^{(T)},Q^{(T)}$ for the components at horizon $T$ and
$L^{(T)}=N-Q^{(T)}$. For $U\leq T$ we have pathwise
\[
 (J^{c,T})^U=J^{c,U}.
\]
On one full-measure set, at every time,
\[
 (P^{(T)})^U=(P^{(U)})^U,\qquad
 (Q^{(T)})^U=(Q^{(U)})^U,\qquad
 (L^{(T)})^U=(L^{(U)})^U.
\]
The localizing sequences used within the individual horizons need not agree.
\end{lemma}
\begin{proof}
The large-jump set in $(0,U]$ is contained in that in $(0,T]$.
Any time in the difference exceeds $U$ and does not contribute to the
finite sum at $t\wedge U$. This proves the first identity. Next consider
\[
 A=(P^{(T)})^U-(P^{(U)})^U
   =(Q^{(U)})^U-(Q^{(T)})^U.
\]
The first expression is predictable and has finite variation over the
whole time axis, because of the finite horizon stop. The second is a
difference of c\`adl\`ag local martingales. Predictable finite-variation
rigidity makes $A$ indistinguishable from its initial value, which is zero.
This proves the identity for $P$, then for $Q$, and finally for $L=N-Q$.
When $U=0$, the identities follow directly from zero initial values.
\end{proof}
Construct these decompositions for $T_n=n+1$, always starting from the
same original local martingale. A countable intersection gives all ordered
overlaps, for all times and all $m\leq n$, on one full-measure set.
The jump bounds of the same $P^{(T_n)}$ can also be made simultaneous in $n$.

\begin{lemma}[Gluing the predictable component over all time]
The coherent horizon family gives a zero-start predictable,
right-continuous locally finite-variation process $P$ such that, on one
full-measure set, for every $n,t$,
$P_{t\wedge T_n}=P^{(T_n)}_{t\wedge T_n}$.
On one full-measure set $|\Delta P_t|\leq c$ for every $t$.
\end{lemma}
\begin{proof}
Stop $P^{(T_n)}$ at $T_n$ to obtain coordinates of global finite variation.
Apply the appropriate ordered overlap according to the order of the two
indices: the coordinates agree after stopping at $T_n\wedge T_m$.
The deterministic schedule is increasing and exhaustive, so predictable
finite-variation gluing applies. Its initial value is almost surely zero.
The usual conditions make the exceptional set measurable at time zero;
set the whole path to zero on that set.
Take a countable intersection to preserve every horizon identity at once.
For a given $t$, choose $n$ with $t\leq T_n$. The left-jump identity before
and at the closed stop gives
$|\Delta P_t|=|\Delta P^{(T_n)}_t|\leq c$.
\end{proof}

\begin{lemma}[Common refinement of local-martingale coordinates]
Suppose zero-start c\`adl\`ag local martingales $X^n$ satisfy
$(X^n)^{\tau_n\wedge\tau_m}\sim(X^m)^{\tau_n\wedge\tau_m}$ for an
almost-surely increasing exhaustive stopping sequence $(\tau_n)$.
There is a c\`adl\`ag local martingale $Y$ such that
$Y^{\tau_n}\sim(X^n)^{\tau_n}$ for every $n$.
\end{lemma}
\begin{proof}
Let $(\eta^n_j)_j$ localize $X^n$. Choose $k(n)$ by a countable extraction
so that $\sigma_n=\tau_n\wedge\eta^n_{k(n)}$ tends to infinity almost surely.
Put $\rho_n=\inf_{j\geq n}\sigma_j$. Right continuity of the filtration
makes it a stopping time; the sequence increases to infinity and satisfies
$\rho_n\leq\tau_n\wedge\eta^n_{k(n)}$.
Zero initial values remove the time-zero indicator in a localizing stop,
and optional stopping makes $(X^n)^{\rho_n}$ a true martingale.
If $W$ is the pointwise limit of the original compatible family, then
$W^{\rho_n}\sim(X^n)^{\rho_n}$. The refined coordinates are therefore
compatible, and true-martingale gluing applies. Its result $Y$ agrees with
$W$ at all $\rho_n$, hence $Y\sim W$ by exhaustion. Stopping again at the
original $\tau_n$ proves the required identities.
\end{proof}

\begin{proposition}[A global Doléans--Dade--Yen decomposition]
For a zero-start c\`adl\`ag local martingale $N$ and $c>0$, there are
zero-start c\`adl\`ag local martingales $L,Q$ with $N=L+Q$, with $Q$ locally
of finite variation, and with $|\Delta L_t|\leq2c$ for every time on one
full-measure set.
\end{proposition}
\begin{proof}
Apply common refinement and gluing to $Q^{(T_n)}$ from the same horizon
family. The horizon identities give local bounded variation and zero
initial value almost surely. Correct the paths to zero on the exceptional
time-zero measurable set, making both properties hold everywhere.
The processes $P,Q$ glued from that family satisfy
$(J^{c,T_n})^{T_n}=(Q+P)^{T_n}$ simultaneously over all horizons.
Set $L=N-Q$. For $t\leq T_n$,
$\Delta L_t=\Delta(N-J^{c,T_n})_t+\Delta P_t$.
For any positive $t$ choose a horizon containing it and add the two bounds
$c$. The time-zero jump is zero. The other properties of $L$ follow by
subtracting local martingales.
\end{proof}
Apply the quadratic construction below to this same global $L$.
Its application to a finite-horizon $L^T$ is another instance of that construction.

\begin{lemma}[Quadratic localization of a bounded-jump local martingale]
Let $M$ be a zero-start c\`adl\`ag local martingale, with deterministic
jump bounds $|\Delta M_t|\leq b_n$ on each integer horizon $n+1$, where
$b_n\geq0$. Let $(\eta_n)$ be a local-martingale localizing sequence and put
\[
 \theta_n=\inf\{t\geq0:|M_t|>n+1\},\qquad
 \rho_n=\eta_n\wedge\theta_n\wedge(n+1).
\]
This sequence is exhaustive, each $M^{\rho_n}$ is a true martingale, and
simultaneously at all times
\[
 |M^{\rho_n}_t|\leq n+1+b_n\quad\text{a.s.}
\]
In particular, each coordinate's terminal value belongs to $L^2$.
\end{lemma}
\begin{proof}
Both the horizon-capped absolute passages and $(\eta_n)$ are localizing
sequences, as is their minimum. At zero, the zero initial value identifies
the indicator-corrected localizing stop with the usual closed stop, even
on paths with $\eta_n=0$. Thus $M^{\eta_n}$ is a true martingale.
Bounded optional stopping at $\rho_n\leq\eta_n$ gives the same property
for $M^{\rho_n}$. Before passage the left-limit magnitude is at most
$n+1$, and the boundary jump adds at most $b_n$. This proves the all-time
bound. The deterministic majorant is in $L^2$, so the terminal value is
square integrable.
\end{proof}
Apply the quadratic construction on each finite horizon and glue at the
minimum of the coordinate stops, obtaining one global optional
quadratic variation $V$. The localizer root of this same $V$ equals the
coordinate's terminal root almost surely, so their expectation costs agree.
Applying Davis's inequality to each coordinate gives
\[
 E\sup_{t\leq n+1}|M^{\rho_n}_t|\leq6E\sqrt{V_{\rho_n}}.
\]
The supremum is detected on a countable time set by c\`adl\`ag regularity.

For the fixed-horizon decomposition take $M=L^T$ and $b_n=2c$.
Thus the quadratic and Davis estimates apply to the same chosen DDY component.
This $V$ is the quadratic variation of $L^T$, not of the original $N$.
The coordinate estimates do not assert global $H^1$ integrability of
unstopped $L^T$.
